\section{Related Work}
\label{sec:related}

\subsection{Convolutional Segmentation Architectures}

The formulation of segmentation as dense per-pixel classification by fully convolutional networks \cite{ref16} established the basis upon which subsequent architectures were constructed. The U-Net \cite{ref4} adapted this formulation to biomedical imaging by pairing a contracting encoder with a symmetric expanding decoder and joining the two by skip connections, so that the fine spatial detail lost to successive downsampling could be reintroduced during reconstruction. This design proved well suited to the small annotated datasets characteristic of medical imaging and remains the reference point for the field. Later work observed that the direct concatenation of encoder and decoder features couples representations of differing semantic level, and UNet++ \cite{ref6} accordingly interposed a nested arrangement of dense convolutional blocks along the skip pathways in order to reduce this disparity. These architectures recover the lesion region reliably, but the locality of the convolution operator restricts the spatial extent over which evidence can be aggregated, which is a limitation where the margin must be inferred from context beyond its immediate neighbourhood.

\subsection{Attention Mechanisms}

Attention mechanisms were introduced into this setting as a means of weighting features according to their informativeness rather than treating them uniformly. Squeeze-and-excitation blocks \cite{ref17} recalibrate feature channels by a globally pooled descriptor, allowing informative channels to be emphasised. The Attention U-Net \cite{ref5} applies a related principle spatially, gating each skip connection so that regions irrelevant to the target are suppressed before the encoder features are passed to the decoder. Such mechanisms improve the selectivity of the representation, and the gating principle they establish --- that the contribution of a feature should be determined from the feature itself rather than fixed in advance --- is directly relevant to the combination of heterogeneous representations. They do not, however, alter the receptive field of the underlying convolutions, and the constraint on long-range aggregation persists.

\subsection{Transformer-Based Segmentation}

The vision transformer \cite{ref7} addressed this constraint by dispensing with convolution in favour of self-attention over image patches, admitting global interactions at every layer. Because pure transformer encoders forgo the inductive biases that make convolutional networks efficient on limited data, hybrid designs were proposed for medical segmentation. TransUNet \cite{ref8} employs a convolutional encoder followed by transformer layers and retains a convolutional decoder, thereby combining local feature extraction with global context modelling, while Swin-UNet \cite{ref18} constructs an encoder--decoder entirely from hierarchical windowed attention blocks. These architectures capture context effectively, but they carry substantially greater computational and data requirements, which is a material consideration given the size of publicly available dermoscopic collections.

\subsection{Foundation Models for Medical Segmentation}

The difficulty of training large models on limited annotated data has motivated the adaptation of visual foundation models pretrained on very large and diverse collections. The Segment Anything Model \cite{ref9} provides a promptable architecture whose image encoder yields broadly transferable representations, and its successor \cite{ref10} extends the family with a hierarchical encoder and improved efficiency. Direct application to medical images is nevertheless impaired by the domain gap between natural and medical imagery, and adaptation is required; MedSAM \cite{ref11} addresses this by fine-tuning on a large assembled corpus of medical images. Fine-tuning an encoder of this scale is costly and, on small datasets, carries a risk of overfitting, which has led to approaches that retain the encoder in a frozen state and confine the trainable capacity to modules operating on its outputs. Under such a scheme the adaptation mechanism, rather than the backbone, governs how faithfully fine structure is preserved.

\subsection{Frequency-Domain Learning}

A complementary line of work processes features in the frequency domain rather than the spatial domain alone. Fast Fourier convolution \cite{ref13} incorporates a spectral branch within the convolutional block so that a global receptive field is obtained at modest cost, and global filter networks \cite{ref14} replace self-attention outright with learnable filters applied to the Fourier transform of the feature map. The interest of these methods for boundary delineation lies in the correspondence between spectral content and spatial structure: abrupt intensity transitions of the kind associated with object margins are represented predominantly in the higher-frequency components, whereas the smoothly varying appearance of homogeneous regions is concentrated in the lower-frequency components. Operating in this domain therefore affords a means of addressing margin-related content more directly than is possible within a single spatial representation.

\subsection{Boundary-Aware Segmentation}

Recognition that region-overlap objectives exert limited pressure upon the contour has prompted the development of boundary-aware formulations. Boundary loss \cite{ref19} supplements a region objective with a distance-weighted term computed from the ground-truth contour, and related work has sought to reduce contour distances directly during optimisation. A parallel concern arises in evaluation: region-overlap criteria are area-based and are correspondingly insensitive to modest displacements of the contour, so that methods differing appreciably in margin fidelity may be assigned similar scores. The Boundary IoU measure \cite{ref12} was proposed in response, restricting the overlap computation to a band about the contour so that boundary quality is assessed separately from region extent. Distance-based measures such as the Hausdorff distance and the average symmetric surface distance serve a comparable purpose in medical image analysis.

\subsection{Summary and Research Gap}

These lines of work converge upon the framework from which the present study proceeds. The Original MAFF-Net \cite{ref15} adapts a frozen foundation backbone to dermoscopic segmentation through multi-scale perception adaptation and frequency-guided feature fusion, together with an auxiliary reconstruction branch, and demonstrates that spatial and spectral information can be integrated productively within a single framework. Three observations nevertheless remain. First, in that design and in comparable approaches the spectral and spatial content are combined within one fusion route rather than being maintained as separate streams; representations of markedly different statistical character are thereby merged without prior reconciliation, and the low-frequency content describing the lesion interior may predominate. Second, where two heterogeneous representations are combined, the weighting between them is commonly fixed rather than determined from the features themselves, notwithstanding that the gating principle established in the attention literature is directly applicable. Third, and cutting across the foregoing, methods intended to improve margin delineation are still assessed principally by region-overlap criteria, which, for the reasons set out above, are ill-suited to detecting the very property such methods are designed to improve.

The present work addresses these three points together. It retains the frozen foundation backbone and the established decoding pathway, separates the spatial and spectral representations into parallel streams, reconciles their statistics before combination, and merges them by a learned gate whose weighting varies with the features presented to it. It further evaluates the resulting model with region-based and boundary-based measures in combination, so that the effect upon the contour may be observed independently of the effect upon region extent.
